%subfiles使用時はdocument環境内にサブファイルも記入する
%サブファイル内のプリアンブルはmain.texのビルド時は無視される
\documentclass[../../main.tex]{subfiles} %[]内にこのドキュメントからmain.texまでのpathを記入

\graphicspath{{../../fig/}} %このドキュメントから画像ディレクトリまでのpathを記述
\setcounter{chapter}{0} %チャプター番号の調整

%***************************************** ＜Main body＞ ********************************************%
\begin{document}

\setcounter{page}{0} %ページ番号をゼロに設定
\pagenumbering{arabic} %本文用のページ番号付け

\chapter{提案手法}\label{chap:introduction}


\section{Human–Robot Interaction with Correction and Suggestion (HRI-CAS)}

\textbf{Human–Robot Interaction with Correction and Suggestion (HRI-CAS)} 
addresses several limitations observed in conventional MR-based HRC systems, 
including insufficient utilization of environmental information, 
lack of stable posture guarantees, 
and inadequate autonomy in base relocation strategies.

This system is built upon the foundation of 
\textit{Immersive Robot Teleoperation Based on User Gestures in Mixed Reality Space (IRT-MRO)} 
and improves upon the issues that remained in IRT-MRO.  
Whereas IRT-MRO employed a Unity-based holographic application on the HoloLens 2, 
the proposed system adopts a Unity-based passthrough application using the Meta Quest.  
Although both devices are capable of projecting real-world objects into MR space, 
they differ significantly in their display methods.  

The HoloLens 2 adopts an optical see-through display, 
allowing the user to perceive both the real environment and virtual objects through a transparent screen.  
However, its limited field of view and restricted brightness have been widely noted.  
In contrast, as illustrated in Figure~1, the Meta Quest employs a passthrough display system, 
which projects images captured by external cameras onto the display, 
enabling a wide field of view and high-resolution rendering.  
This provides the user with a more natural and immersive MR experience, 
enhancing both environmental understanding and intuitive operation.
% Insert figure

\subsection{Implementation Platform and System}
In this study, the experimental platform adopted was the \textbf{KUKA YouBot}.  
The robot is a mobile manipulator equipped with an omnidirectional base, 
a 5-DOF manipulator, and a two-finger gripper.  
Within the proposed system, changes in the position and orientation of target objects 
detected in MR space, as well as user manipulations, 
are reflected directly in the robot’s motion.  
This ensures smooth and seamless interaction between the human and the robot.  
Furthermore, the robot’s behavior is immediately fed back into MR space, 
providing the operator with precise visual information in real time.

The graphical architecture of the system is illustrated in Figure~2 
and is composed of the following three modules:
% Insert figure

\medskip
\noindent
\textbf{Target Bottle Inference (TBI) module}:  
This module tracks the user’s gaze, head orientation, and hand motion within MR space, 
extracting geometric relationships with candidate bottles as features.  
These multimodal features are then processed through an integrated scoring framework 
to infer the bottle that the user intends to manipulate.  
By jointly evaluating diverse behavioral cues, 
the TBI module enables highly accurate and real-time target selection.

\medskip
\noindent
\textbf{Base Relocation based on Inverse Reachability Map (BRI) module}:  
When the target bottle lies outside the robot’s reachable workspace, 
this module becomes active.  
First, an Inverse Reachability Map (IRM) is generated based on the Reachability Map (RM), 
and candidate base positions that make the target reachable are extracted.  
Next, an optimization problem is solved considering both the IRM score 
and the movement cost, yielding the optimal base configuration.  
Additionally, candidate paths are visualized within MR space, 
allowing users to refine or adjust the proposed routes interactively.

\medskip
\noindent
\textbf{Human-Intent Adaptive Grasp Manipulation (HIGM) module}:  
This module aims to reconcile the differences between the user’s intended manipulation 
and the robot’s physical and kinematic constraints.  
Specifically, it integrates reference postures obtained from hand-tracking information 
with robot-preferred pregrasp postures, 
and applies nullspace optimization to adjust joint configurations.  
As a result, the user can achieve stable grasping with natural operational feedback, 
without needing to be aware of the robot’s inherent kinematic constraints.

\medskip
Through the integration of these modules, 
the proposed system achieves adaptive and intuitive motion planning 
that was not sufficiently addressed in IRT-MRO.  
The robot is now able to select actions according to the user’s intended target, 
while providing adaptive grasp correction without requiring the user to 
account for robot-specific constraints.  
This significantly reduces dependency on the user’s operational expertise 
and ensures smooth, stable operation even for novice users.

\subsection{MR System Implementation}
The MR system is implemented on \textbf{Unity} and \textbf{MRTK3 (Mixed Reality Toolkit)}, 
in compliance with the OpenXR standard.  
The user can manipulate MR objects directly through hand gestures, 
offering an intuitive and natural input interface.  
Moreover, by presenting the robot’s behavioral intentions through MR objects, 
the system improves the operator’s understanding of the robot’s decision-making process, 
thereby enhancing human–robot collaboration.




\end{document}